\newsection{Reguläre Sprachen und Automaten}{task-2}

\newcommand{\statement}[1]{
  \item{
    \noindent\begin{minipage}[t]{0.65\linewidth} % Left side (text)
      #1
    \end{minipage}
    \hfill
    \begin{minipage}[t]{0.31\linewidth} % Right side
      \vfill % Used to center vertically
      \begin{center}
        \setlength{\tabcolsep}{10pt} % Sets the column padding
        \begin{tabular}{@{}c c@{}}
          Wahr & Falsch \\
          $\square$ & $\square$
        \end{tabular}
      \end{center}
      \vfill % Used to center vertically
    \end{minipage}

    \vspace{10pt}
    \begin{minipage}{\linewidth}
    Begründung: \hrulefill \\[10pt] \hrule \vspace{20pt} \hrule
    \end{minipage}
  }
}

\begin{enumerate}
  \item {
    Geben Sie für den folgenden regulären Ausdruck die \textbf{Ableitungsregeln}
    einer Grammatik an.
    Verwenden Sie dabei so \textbf{wenige Regeln} wie möglich.

    \[
      R = b(aa + b)^*
    \]

    \vspace{10pt}

    Ableitungsregeln:
  }
  \vspace{\fill}
  \item {
    Zeichnen Sie für die folgende Sprache einen Deterministischen Endlichen
    Automaten (DFA) als Zustandsübergangsdiagramm:

    \[
      L = \{w \in \{0, 1\}^*\mathbin{:} w \text{~beginnt mit 1 und endet auf 0}\}
    \]
  }
  \vspace{\fill}
  \newpage
  \item {
    Geben sind die vier folgenden Aussagen über reguläre Sprachen.
    Entscheiden Sie jeweils, ob diese Aussagen wahr oder falsch sind und geben
    Sie eine Begründung an.

    \begin{enumerate}
      \setlength\itemsep{10pt}
      \statement{
        Wenn es bei einem NFA mindestens einen Übergang gibt, sodass das gelesene
        $w \in \Sigma{^*}$ in einem verwerfenden Zustand landet,
        so akzeptiert der NFA $w$ nicht.
      }
      \statement{
        Ein minimaler NFA hat höchstens so viele Zustände wie ein DFA,
        der die gleiche Sprache entscheidet, mit minimaler Anzahl an Zuständen.
      }
      \statement{
        Sei $k \in \mathbb{N}$ beliebig, aber fest.
        Jede Sprache, die nur Wörter der Länge $k$ akzeptiert, ist regulär.
      }
      \statement{
        Jede nicht reguläre Sprache kann durch das Pumping-Lemma auch als solche
        identifiziert werden.
      }
    \end{enumerate}
  }
  \newpage
  \item {
    Beweisen Sie mit Hilfe des Pumping-Lemmas, dass folgende Sprache über dem
    Alphabet $\Sigma = \{a, b\}$ nicht regulär ist:

    \[
      L = \{w \in \Sigma ~| ~|w|_a = |w|_b + 1\}
    \]
  }
\end{enumerate}